\documentclass[11pt]{article}
\usepackage[a4paper,margin=1in]{geometry}
\usepackage{parskip}
\usepackage{hyperref}
\hypersetup{colorlinks=true, urlcolor=blue, linkcolor=blue}
\pagestyle{empty}

\begin{document}

\begin{flushright}
September 3, 2026
\end{flushright}

To the Editor-in-Chief\\
\textit{Project Management Horizons} (MDPI)

\bigskip
\noindent\textbf{Re: Submission of manuscript ``Evaluating Large Language Models for Automated Risk Register Generation from Project Documentation: A Benchmark Study''}

\bigskip
Dear Editor,

We are pleased to submit our manuscript, ``Evaluating Large Language Models for Automated Risk Register Generation from Project Documentation: A Benchmark Study,'' for consideration for publication in \textit{Project Management Horizons}.

This work is fundamentally an artificial intelligence evaluation study: its subject is the capability and failure behavior of large language models (LLMs) -- proprietary and open-weight generative AI systems -- when applied to a real project management task. Generative AI is increasingly proposed for project management work, yet published evaluation of it has concentrated almost entirely on predictive analytics for schedule, cost, and quality, leaving the generative capability of these AI systems comparatively under-studied. Within the narrower literature that does examine LLMs for project risk work specifically, evaluation has so far relied on expert-opinion panels rather than comparison against real, published, human-authored artifacts, and has been almost exclusively single-model. To our knowledge, this is the first study to benchmark end-to-end risk-register \emph{generation} by AI -- not classification of already-identified risks -- against genuine, published, human-authored risk registers, across three LLMs (one proprietary, two open-weight) and three prompting strategies, using 21 real projects with public planning documentation (18 World Bank Project Appraisal Documents and 3 UK government business cases). Across 302 real evaluation runs, we find that AI-generated registers recover just over half of a human-authored register's real content (mean recall 0.556) while fewer than a quarter of generated risks correspond to a real ground-truth risk (mean precision 0.230), with a specific, non-uniform category failure pattern: political and regulatory risk is disproportionately missed by the models, while schedule risk is generated as a near-default output regardless of project fit. We believe this quantified, failure-mode-level account of a widely proposed but under-evaluated AI application is a timely and substantive contribution to the literature on AI in project management practice.

This manuscript sits squarely within the scope of \textit{Project Management Horizons}: it examines an emerging AI technology's application to a core project management deliverable -- the risk register -- and reports both its practical capability and its systematic limitations, directly relevant to project management researchers and practitioners evaluating whether and how to incorporate LLM-based AI tools into risk management practice. We also release the full corpus manifest, extraction and prompting code, and evaluation pipeline to support reproducibility.

We confirm that neither the manuscript nor any parts of its content are currently under consideration for publication with or published in another journal.

All authors have approved the manuscript and agree with its submission to \textit{Project Management Horizons}.

This manuscript has not been previously submitted, in whole or in part, to this or any other MDPI journal.

Thank you for considering our manuscript. We look forward to your response.

\bigskip
Sincerely,

\bigskip
\noindent Madhu Sudhan (corresponding author)\\
Dayananda Sagar College of Engineering, Bengaluru, Karnataka 560078, India\\
\href{mailto:smadhu6364@gmail.com}{smadhu6364@gmail.com}

\bigskip
\noindent Kruthik Jadhav\\
School of Professional Studies, New York University, New York, NY 10003, USA\\
\href{mailto:kj2578@nyu.edu}{kj2578@nyu.edu}

\end{document}
